Am verificat corectitudinea fragmentarii analizand cele trei reguli fundamentale prezentate la curs:
completitudine, reconstructie, disjunctie.

1. Verificare completitudinii
Presupune ca fiecare tuplu din relatia initiala sa se regaseasca in cel putin unul din fragmentele obtinute

Relatia CLIENT a fost fragmentata in functie de tipul clientului astfel:
Clientii care au atributul client_type setat la valoarea 'F' -> CLIENT_FIZIC
Clientii care au atributul client_type setat la valoarea 'J' -> CLIENT_JURIDIC

Constrangerea CHECK (client_type IN ('F', 'J')) ne asigura ca fiecare client apartine obligatoriu unuia
din cele doua fragmente

2. Verificarea reconstructiei
Permite obtinerea relatiei initiale fara pierdere de informatie din fragmentele rezultate

Pentru fragmentarea orizontala reconstructia se face prin reuniune:

CLIENT = CLIENT_FIZIC UNION ALL CLIENT_JURIDIC;

Este asigurata prin view-uri globale (V_TICKET_AGENT), reunind fragmentele si permitand aplicatiei sa
acceseze datele ca pe o singura relatie logica

Pentru fragmentarea verticala reconstructia se face prin join pe cheia primara comuna. Daca informatiile
despre client sunt impartite in date generale si specifice, relatia poate fi refacuta prin join pe client_id

3. Verificarea disjunctiei
Daca acelasi tuplu nu apare simultan in mai multe fragmente (cu exceptia atributelor cheie necesare
reconstructiei fragmentarii verticale) atunci fragmentarea este disjuncta

Fragmentele CLIENT_FIZIC si CLIENT_JURIDIC sunt disjuncte deoarece un client nu poate avea simultan
client_type = 'F' si client_type = 'J'